\documentclass{article}
\usepackage{graphicx} % Required for inserting images
\usepackage{amsmath}  % Essential for matrices, aligned equations, and dots
\usepackage{amssymb}  % Essential for symbols like \mathbb{R} (the real numbers set)
\usepackage{bm}       % Essential for making bold math symbols (vectors/matrices)

\title{Gamma distribution}
\author{Nguyen Phu Vinh}
\date{April 2026}

\begin{document}

\maketitle

\section{Exponential Distribution}

\textbf{How much time will pass before the first event happens?} \\

Assume independence, constant average rate $\lambda$ (e.g. 3 events per year). \\
Then, expected event per given time is: \\
$E[\textbf{events in t}] = \lambda \times t$ \\
Now look at infinitesimal slice of time $dt$, the expected events during that time $dt$ is: \\
$E[\textbf{events in dt}] = \lambda dt$ \\ \\
Because $dt$ is so microscopic, 2 events happening simultaneously during this interval of time is impossible, therefore the the outcome is strictly binary: either exactly \textbf{1} event happens, or \textbf{0} events happen. \\
Note that for binary variable, the expected outcome of an event equals the probability of that same event \\
\textit{(The expected value of a single fair coin flip is 0.5)}. \\
Therefore, probability of \textbf{1} event in $dt$ is: \\
$P[\textbf{1 event in $dt$}] = \lambda dt$ \\
Examine the complement: \\
$P[\textbf{0 events in $dt$}] = 1 - \lambda dt$ \\
\text{Suppose the microscopic time interval } $dt = \dfrac{t}{n}$, where n is the number of slices: \\
$\implies P\left[ \textbf{0 events in }\dfrac{t}{n}\right] = 1 - \dfrac{\lambda t}{n}$\\
Therefore, knowing that events are \textbf{independent}, the probability that no events happen in all time slices (or in time t) is: \\
$P\left[ \textbf{0 events in }{t}\right] = \left(1 - \dfrac{\lambda t}{n}\right)^n$\\
Go back to infinitesimal world, to make this function continuous, we need limits: \\
$\underset{n\rightarrow \infty}{lim} \left(1 - \dfrac{\lambda t}{n}\right)^n = e^{-\lambda t}$ \\
\textit{(This is definition of the number $e$, and in general, the definition of exponential function)} 
\newpage
Take complement again, the accumulated probability of getting at least 1 event up to time t, or the CDF is: \\
\textit{(To clarify, the complement of zero is at least once)} \\
$F(t) = 1 - e^{-\lambda t}$ \\
Differentiate the CDF to get the PDF: \\
$f(t) = \dfrac{d}{dt}(1 - e^{-\lambda t}) =\lambda e^{-\lambda t}$ \\
The probability of at least 1 event happening in a time interval [a, b] is: \\
$P(a<T<b) = \int_{a}^{b} f(t)dt = \int_{a}^{b}\lambda e^{-\lambda t}dt$

\maketitle

\section{Poisson Distribution}
\textbf{How many events happen in a fixed window of time?} \\

Using the same line of logic up until we have binary outcomes: either exactly \textbf{1} event happens, or \textbf{0} events happen. \\
We have: \\
$P[\textbf{1 event in $dt$}] = p = \lambda dt = \dfrac{\lambda t}{n}$  \\
$P[\textbf{0 events in $dt$}] = 1 - p = 1- \lambda dt = 1 -  \dfrac{\lambda t}{n}$  \\
Now, instead of asking whether event happens at least once or does not happen at all (Exponential Distribution), we ask: what is the probability of getting exactly $k$ events across $n$ slices. \\
\textit{(Note that each slice can carry at most 1 event, that is why Binomial is used in this case)} \\
We use the Binomial Distribution formula: \\
$P(X=k) = \binom{n}{k} p^k(1-p)^{n-k}$ \\

Substitute $p = \dfrac{\lambda t}{n}$ \\
$P(X=k) = \dfrac{n!}{k!(n-k)!} \left( \dfrac{\lambda t}{n} \right)^k\left( 1 - \dfrac{\lambda t}{n} \right)^{n-k}$ \\

Factor $\dfrac{1}{k!}$ out and isolate terms with $n$ (to take limit $n \rightarrow \infty$ later): \\
$P(X=k) = \dfrac{1}{k!} \ \dfrac{n(n-1)...(n-k)...1}{(n-k)(n-k-1)...1} \ \dfrac{(\lambda t)^k}{n^k} \ \left( 1 - \dfrac{\lambda t}{n} \right)^{n}  \left( 1 - \dfrac{\lambda t}{n} \right)^{-k}$ \\


$P(X=k) = \dfrac{n(n-1)...(n-k+1)}{k!} \ \dfrac{(\lambda t)^k}{n^k} \ \left( 1 - \dfrac{\lambda t}{n} \right)^{n}  \left( 1 - \dfrac{\lambda t}{n} \right)^{-k}$ \\

Commute position of $n^k$ and $k!$: \\
$P(X=k) = \dfrac{n(n-1)...(n-k+1)}{n^k} \ \dfrac{(\lambda t)^k}{k!} \ \left( 1 - \dfrac{\lambda t}{n} \right)^{n}  \left( 1 - \dfrac{\lambda t}{n} \right)^{-k}$ \\

Take limit $n \rightarrow \infty$: \\
$P(X=k) = lim_{n\rightarrow \infty} \ \dfrac{n(n-1)...(n-k+1)}{n^k} \ \dfrac{(\lambda t)^k}{k!} \ \left( 1 + \dfrac{-\lambda t}{n} \right)^{n}  \left( 1 + \dfrac{-\lambda t}{n} \right)^{-k}$ \\


\textbf{Examine each term:} \\ 
- Term 1: \\
$$lim_{n\rightarrow \infty} \ \dfrac{n(n-1)...(n-k+1)}{n^k} $$ \\
There are exactly $k$ terms in the numerator, therefore, as $n \rightarrow \infty$, \\ $n(n-1)...(n-k+1)  = n^k$ \\
$\implies lim_{n\rightarrow \infty} \ \dfrac{n(n-1)...(n-k+1)}{n^k} = \dfrac{n^k}{n^k} = 1$ \\
- Term 2:
$$\lim_{n \to \infty} \dfrac{(\lambda t)^k}{k!}$$ \\
The term is constant, not dependent on $n$ \\
$\implies \lim_{n \to \infty} \dfrac{(\lambda t)^k}{k!} = \dfrac{(\lambda t)^k}{k!}$ \\
- Term 3:  \\
$$\lim_{n \to \infty} \left(1 + \dfrac{-\lambda t}{n}\right)^n$$ \\
This is the definition of an exponential function \\
$\implies \lim_{n \to \infty} \left(1 + \dfrac{-\lambda t}{n}\right)^n = e^{-\lambda t}$ \\
- Term 4: \\
$$\lim_{n \to \infty} \left(1 + \dfrac{-\lambda t}{n}\right)^{-k}$$ \\
The term $\dfrac{-\lambda t}{n} \rightarrow 0 $\\
$\implies \lim_{n \to \infty} \left(1 + \dfrac{-\lambda t}{n}\right)^{-k} = (1 +0)^{-k} = 1$ \\

\textbf{Piece all together} \\
$$P(X=k) = \dfrac{(\lambda t)^k e^{-\lambda t}}{k!}$$

\newpage
\maketitle

\section{Gamma Distribution}
\textbf{The continuous waiting time for the $k$-th event.} \\

Let $T$ be the continuous waiting time until the $k$-th event. Then, the probability that waiting time for $k$-th event is GREATER than $t$ is:\\
$$P(T>t)$$ 
Logically, if the $k$-th event does occur after $t$, strictly less than $k$ events must occur before $t$ (could be $0, 1,.., k-1$ events). \\
The probability that less than $k$ events occur up to time $t$ is: \\
$$P(\text{events up to time } t<k)$$ 
Since these two views are identical, we have: \\
$$P(T>t) = P(\text{events up to time } t<k)$$
Since we are talking about probability of some events up to time $t$, it is only rational to think about Poisson distribution.  \\
Recall that probability of having exactly $k$ events up to time $t$ is:
$$P(X=k) = \dfrac{(\lambda t)^k e^{-\lambda t}}{k!}$$
Look at the complement, the probability of NOT having exactly $k$ events up to time $t$ is:  
$$\bar{P}(X=k) = 1- \dfrac{(\lambda t)^k e^{-\lambda t}}{k!}$$
This includes all cases from $X = \{0,1,...,k-1,k,...,\infty\}$ \\
However, we only desire $X<k$ cases. \\
Let us look at a complement rule example: \\
$P(X=0) + \bar{P}(X=0)  = 1$ \\
The complement includes ALL cases that is not EXACTLY 0 events. \\
$P(X=0) + [P(X=1) + P(X=2) + ...+P(X=\infty)] = 1 $\\
Pull all the desired cases out of the sum: \\
$[P(X=0) + P(X=1) +...+ P(X=k-1)] +[P(X=k)+ ...+P(X=\infty)] = 1$\\
Simplify to: \\
$$P(X<k)= 1 - P(X\geq k)$$
Substitute with Poisson function: \\
$$P(X < k) = \sum_{i=0}^{k-1} \dfrac{(\lambda t)^i e^{-\lambda t}}{i!}$$
Substitute with the initial premise: $P(T>t) = P(\text{events up to time } t<k)$ \\ and to turn this into a function of time $t$ \\
$$P(T>t) = \sum_{i=0}^{k-1} \dfrac{(\lambda t)^i e^{-\lambda t}}{i!}$$
$\implies \textbf{The probability that waiting time for $k$-th event is GREATER than $t$}$ \\ \\
\text{The complement is:} \\
$$P(T \leq t) = 1 - \sum_{i=0}^{k-1} \dfrac{(\lambda t)^i e^{-\lambda t}}{i!}$$ \\
$\implies \textbf{The probability that waiting time for $k$-th event is WITHIN $t$} $ \\
\textit{Note: $P(T \leq t)$ is the CDF} \\ \\
\textbf{Building the PDF} \\
Let the CDF be $F(t)$, the derivative w.r.t $t$ is:
$$\dfrac{dF(t)}{dt} =\dfrac{d}{dt} \left[1- \sum_{i=0}^{k-1} \dfrac{(\lambda t)^i e^{-\lambda t}}{i!}\right]$$
Simplify term $(\lambda t)^i = \lambda^i t^i$:  \\
$$\dfrac{dF(t)}{dt} =\dfrac{d}{dt} \left[1- \sum_{i=0}^{k-1} \dfrac{\lambda^i t^i e^{-\lambda t}}{i!}\right]$$

Use product rule: $\dfrac{d}{dt}[u(t)v(t)] = u'(t)v(t) + u(t)v'(t)$

$$\dfrac{dF(t)}{dt} =- \sum_{i=0}^{k-1} \dfrac{i\lambda^i t^{i-1} e^{-\lambda t} + (-\lambda)\lambda^i t^i e^{-\lambda t}}{i!}$$
$$\dfrac{dF(t)}{dt} =- \sum_{i=0}^{k-1} \left( \dfrac{\lambda^i t^{i-1} e^{-\lambda t}}{(i-1)!} - \dfrac{\lambda^{i+1} t^i e^{-\lambda t}}{i!} \right)$$

\newpage

Expand the sum: \\
+) $i=0$ 
$$0 - \underbrace{\dfrac{\lambda^1 t^0 e^{-\lambda t}}{0!}} $$
+) $i=1$ 
$$\underbrace{ \dfrac{\lambda^1 t^0 e^{-\lambda t}}{0!}} - \dfrac{\lambda^2 t^1 e^{-\lambda t}}{1!}$$
+) $i=2$ 
$$\dfrac{\lambda^2 t^1 e^{-\lambda t}}{1!} - \dfrac{\lambda^3 t^2 e^{-\lambda t}}{2!}$$
+) $i=k-1$
$$\dfrac{\lambda^{k-1} t^{k-2} e^{-\lambda t}}{(k-2)!} - \dfrac{\lambda^k t^{k-1} e^{-\lambda t}}{(k-1)!}$$

\textbf{Look at underbraced term, notice that the right term of $i=0$ equals to the left term of $i = 1$, this pattern persists throughout the entire sum} \\
$\implies \text{The only remaining term is:}$
$$\dfrac{dF(t)}{dt} = -\left(- \dfrac{\lambda^k t^{k-1} e^{-\lambda t}}{(k-1)!} \right)$$
The PDF is:
$$f(t) =  \dfrac{\lambda^k t^{k-1} e^{-\lambda t}}{(k-1)!} $$
The area under the curve:
$$F(t) = \int_0^t \dfrac{\lambda^k x^{k-1} e^{-\lambda x}}{(k-1)!} dx$$


\end{document}
